%!TEX root=../../note.tex
\subsection{Basic Set Notation and Operations}
Set notation describes which objects belong to a collection. Union,
intersection, and difference combine sets by applying logical conditions
to membership.

\begin{definition}[Empty Set and Cardinality]
A \emph{set} is a collection of distinct elements. The \emph{empty set} $\emptyset$ has no elements. For a finite set $S$, its \emph{cardinality} $\abs{S}$ is its number of elements.
\end{definition}
We can describe a set by listing its elements, selecting elements with a property, or generating values:
\begin{itemize}
    \item Listing: $\set{2,3,4,6}$.
    \item Selection: $\set{n\in\Z : 0<n<10 \text{ and } n \text{ is even}}=\set{2,4,6,8}$.
    \item Generation: $\set{2k : k\in\set{1,2,3,4}}=\set{2,4,6,8}$.
\end{itemize}
\begin{example}
Order and repetition do not change a set: $\set{1,2,2}=\set{2,1}$. Also, $\abs{\emptyset}=0$, while $\abs{\set{\emptyset}}=1$: the latter set contains the empty set as its one element.
\end{example}

\subsubsection{Set-Builder Notation}
\begin{definition}[Set-Builder Notation]
The \emph{selection form} $\set{x\in U : P(x)}$ contains exactly the elements of $U$ for which $P(x)$ is true. The \emph{generation form} $\set{f(x) : x\in U}$ contains the values $f(x)$ for $x\in U$; its elements are these outputs, not the inputs $x$.
\end{definition}
\begin{example}
    $\set{x^2 : x\in\set{-2,-1,0,1,2}}=\set{0,1,4}$: different inputs may give the same output, which appears only once.
\end{example}
We can combine the forms: $\set{f(x) : x\in U,\ P(x)}$ first selects the inputs satisfying $P$, then forms their outputs. For example,
\[
    \set{3a : a\in\Z,\ 2\mid a}=\set{6k : k\in\Z}.
\]

\subsubsection{Intervals}
For real numbers $a<b$, square brackets include an endpoint and round brackets exclude it:
\begin{gather*}
    [a,b]=\set{x\in\R : a\leq x\leq b}, \qquad
    (a,b)=\set{x\in\R : a<x<b},\\
    [a,b)=\set{x\in\R : a\leq x<b}, \qquad
    (a,b]=\set{x\in\R : a<x\leq b}.
\end{gather*}

\subsubsection{Set Operations}
\begin{definition}[Set Operations]
The \emph{union} $S\cup T$ collects elements in either set; the \emph{intersection} $S\cap T$ keeps elements in both. The \emph{difference} $S-T$ keeps elements in $S$ but not $T$:
\begin{gather*}
    S\cup T=\set{x : x\in S\text{ or }x\in T}, \qquad
    S\cap T=\set{x : x\in S\text{ and }x\in T},\\
    S-T=\set{x : x\in S\text{ and }x\notin T}.
\end{gather*}
When $S$ is a subset of a universal set $U$, its \emph{complement} is $\overline{S}=U-S$.
\end{definition}

\mdfsetup{nobreak=true}
\begin{example}
Let $A=[-1,2)$ and $B=(0,4]$. Then
\begin{gather*}
    A\cup B=[-1,4], \qquad A\cap B=(0,2),\\
    A-B=[-1,0], \qquad B-A=[2,4].
\end{gather*}
The point $2$ belongs to $B$ but not $A$, so it is excluded from the intersection and included in $B-A$.
\end{example}
\mdfsetup{nobreak=false}

For a set $S$ contained in the universal set $U$, the following identities
hold, with complements taken relative to $U$:
\begin{gather*}
    S\cup\emptyset=S, \qquad S\cap\emptyset=\emptyset,\\
    S\cup U=U, \qquad S\cap U=S,\\
    S\cup S=S, \qquad S\cap S=S,\\
    S\cup\overline{S}=U, \qquad S\cap\overline{S}=\emptyset.
\end{gather*}

\subsubsection*{In practice}
\begin{itemize}
    \item \textbf{Reading set-builder notation.} Before the colon: what the elements look like. After it: where the variable ranges and which conditions it must satisfy. To test membership of $y$ in $\set{f(x) : x\in U}$, look for some $x\in U$ with $f(x)=y$.
    \item \textbf{Operations on intervals.} Draw both intervals on a number line and check each endpoint separately: it belongs to the result exactly when the defining condition ($\text{or}$, $\text{and}$, $\text{and not}$) holds for it.
\end{itemize}
